Concretely, \ruleName{eval-binop} evaluates both operands and binds the result of the second:
\begin{mathpar}
\inferrule*[lab={\ruleName{eval-binop}}]
{
  \rho, e \Rightarrow \val{v} \\
  \rho, e' \Rightarrow r
}
{
  \rho, \exBinOp{e}{\odot}{e'} \Rightarrow \mlet{v'}{r}{\hat{\odot}(v, v')}
}

\end{mathpar}
The meaning $\hat{\odot}(v, v')$ of the operator is likewise a partial map from operands to results, with
an abort if Python raises an exception; an expression whose operator meaning is undefined has no
derivation. An abort from the first operand is propagated by a separate rule:
\begin{mathpar}
\inferrule*[lab={\ruleName{eval-binop-fail}}]
{
  \rho, e \Rightarrow \aborts{\kappa}
}
{
  \rho, \exBinOp{e}{\odot}{e'} \Rightarrow \aborts{\kappa}
}

\end{mathpar}
An abort may also originate in a rule rather than be propagated from a premise. Calling a value that is
neither closure nor primitive aborts with $\errType$, because Python raises \pyinline{TypeError}:
\begin{mathpar}
\inferrule*[lab={\ruleName{eval-call-nonfun}}]
{
  \rho, e \Rightarrow \val{v} \\
  \rho, \seq{e} \Rightarrow \val{\seq{u}} \\
  v \text{ neither closure nor primitive}
}
{
  \rho, \exCall{e}{\seq{e}} \Rightarrow \aborts{\errType}
}

\end{mathpar}
The remaining rules are of the same form and are given in full in \specification.
